\documentclass{article}
\usepackage{amsmath, amssymb, graphicx}

\title{Status Report: Lanzarini Optimizer Series (V18--Final Benchmark)}
\author{Experimental Optimization Study}
\date{}

\begin{document}

\maketitle

\section{Overview}

This document summarizes the current state of the Lanzarini Optimizer research, including all experimental stages from early variants (V18--V21) to the final benchmark comparison against standard optimization algorithms.

The goal of this report is to provide a clear, reproducible and structured description of:
\begin{itemize}
    \item optimization formulation
    \item experimental setup
    \item observed results
    \item current interpretation of findings
\end{itemize}

\section{Core Optimization Framework}

We consider the standard optimization problem:
\[
\min_{\theta} \mathcal{L}(\theta)
\]

The general update rule used in the Lanzarini framework is defined as:

\begin{equation}
\Delta \theta_t = \alpha \cdot u_t^{(Adam)} + (1 - \alpha) \cdot u_t^{(grad)}
\end{equation}

where:

\begin{itemize}
    \item $u_t^{(Adam)}$ is the Adam-style adaptive direction
    \item $u_t^{(grad)}$ is the raw gradient direction
    \item $\alpha \in [0,1]$ is a mixing coefficient
\end{itemize}

In later versions, $\alpha$ was either:
\begin{itemize}
    \item fixed (V19--V20)
    \item learned via policy (V21)
    \item empirically optimized (final benchmark)
\end{itemize}

\section{Experimental Setup}

All experiments were conducted under controlled conditions:

\begin{itemize}
    \item Multiple random seeds (5 per method)
    \item Quadratic synthetic objective (initial phase)
    \item CIFAR-style benchmark (final phase)
    \item Fixed learning rate across methods
    \item Same initialization distribution
\end{itemize}

Methods compared:

\begin{itemize}
    \item SGD with momentum
    \item Adagrad
    \item RMSProp
    \item Adam
    \item AdamW
    \item Lanzarini (adaptive interpolation)
\end{itemize}

\section{Quadratic Benchmark Results}

Initial experiments were conducted on a quadratic convex optimization problem:

\begin{itemize}
    \item Lanzarini showed faster convergence than Adam in terms of iterations
    \item Reduced number of steps to reach target loss
    \item Improved stability compared to unconstrained interpolation variants (V18)
\end{itemize}

Key observation:
\begin{equation}
\text{Performance gain correlates strongly with interpolation between gradient directions}
\end{equation}

\section{4-Way Ablation Study}

Ablation results comparing different interpolation strategies:

\begin{itemize}
    \item Adam baseline
    \item Fixed interpolation
    \item Random interpolation
    \item V21 learned policy
\end{itemize}

Result pattern:

\begin{itemize}
    \item All interpolation methods significantly outperform SGD
    \item Differences between fixed, random, and learned policies are smaller than expected
    \item Learned policy (V21) does not consistently dominate simpler alternatives
\end{itemize}

\section{Oracle and Decay Tests}

Additional experiments:

\begin{itemize}
    \item Oracle optimal constant $\alpha$
    \item Linear, cosine, and step decay schedules
\end{itemize}

Observation:

\begin{itemize}
    \item A near-optimal constant $\alpha$ achieves performance close to or better than learned policy
    \item Simple scheduling strategies can match or exceed learned control in some settings
\end{itemize}

\section{Final Benchmark (Realistic Setup)}

Final evaluation on extended benchmark with multiple optimizers:

\begin{verbatim}
Methods evaluated:
- SGD
- Adagrad
- RMSProp
- Adam
- AdamW
- Lanzarini (alpha interpolation)
\end{verbatim}

Measured metrics:
\begin{itemize}
    \item Final loss
    \item Standard deviation across seeds
    \item Execution time
    \item Composite score (loss × time)
\end{itemize}

\section{Key Results (Observed)}

\begin{itemize}
    \item Lanzarini achieves the lowest loss among tested methods
    \item RMSProp is second best in stability-performance tradeoff
    \item Adam and AdamW show similar performance in this setup
    \item SGD and Adagrad perform significantly worse
\end{itemize}

\section{Current Interpretation}

Based on all experiments, the following hypotheses are under evaluation:

\subsection{H1: Learned control advantage}
Performance improvement comes from adaptive policy selection.

\subsection{H2: Interpolation effect}
Performance improvement is primarily due to convex combination of two update directions.

\subsection{H3: Variance reduction}
Performance improvement is caused by reduction in update variance rather than information content.

Current evidence suggests:

\begin{itemize}
    \item Strong contribution from interpolation structure
    \item Weak additional gain from learned policy in current experiments
    \item Significant reduction in variance of updates correlates with performance gain
\end{itemize}

\section{Conclusion}

The Lanzarini optimizer series demonstrates that:

\begin{itemize}
    \item Combining gradient-based and adaptive directions improves convergence stability
    \item The structure of interpolation plays a dominant role in performance gains
    \item Learned control policies provide limited additional benefit in current experimental conditions
\end{itemize}

Further work is required to:
\begin{itemize}
    \item test on larger-scale deep learning tasks
    \item isolate variance reduction effects from true adaptive behavior
    \item evaluate generalization across architectures
\end{itemize}

\end{document}
